% !TeX program = pdflatex
\documentclass[11pt,a4paper]{article}
\usepackage{../../recursos/latex/estilo-notas}

\begin{document}
\cabecalhoaula{04}{Métodos Não Paramétricos (KNN)}{AME §4.1--4.3; ISLP §3.5 e §7.1--7.5}

%==============================================================================
\section*{Objetivos da aula}
%==============================================================================
Ao final desta aula o aluno deve ser capaz de:
\begin{itemize}[leftmargin=1.4cm]
  \item explicar o que distingue um método \emph{não paramétrico} de um
        paramétrico e por que ele troca viés por variância;
  \item descrever a estratégia de \emph{bases} (séries ortogonais, \emph{splines})
        para flexibilizar a regressão linear;
  \item definir o estimador de \textbf{$k$ vizinhos mais próximos} (KNN) e o papel
        de $k$ no balanço viés--variância;
  \item enxergar todos esses métodos como \emph{médias locais} / \emph{suavizadores
        lineares}.
\end{itemize}

\begin{emsala}
Mudança de paradigma em relação às Aulas~01--02: lá fixávamos a forma de $r$
(linear) e estimávamos poucos parâmetros. Aqui deixamos \emph{os dados ditarem a
forma}. Pergunta: \emph{``e se eu não quiser supor forma alguma para
$r(\x)$?''}
\end{emsala}

%==============================================================================
\section{O que é um método não paramétrico?}
%==============================================================================
Um método \textbf{paramétrico} supõe que $r(\x)=\E[Y\mid\X=\x]$ pertence a uma
família indexada por um número \emph{fixo e finito} de parâmetros (ex.: regressão
linear, $p+1$ coeficientes). Um método \textbf{não paramétrico} não impõe uma
forma funcional fixa: a complexidade efetiva do modelo \emph{cresce com $n$}. A
suposição é mais fraca --- tipicamente apenas que $r$ é \emph{suave}.

\begin{ideia}
Em relação aos paramétricos, métodos não paramétricos \textbf{reduzem o viés}
(capturam relações arbitrárias) ao custo de \textbf{aumentar a variância}.
Continuamos no mundo da Aula~01: o segredo é controlar a flexibilidade por um
\emph{tuning parameter} (o $k$ do KNN, o número de nós de um \emph{spline}),
escolhido por
validação cruzada (Aula~03).
\end{ideia}

\begin{observacao}
Há uma ponte importante com a Aula~02: penalizar (Ridge, Lasso) e usar um método não
paramétrico buscam o mesmo ponto do balanço viés--variância, atacando parcelas
opostas. A penalização é para o modelo flexível demais para a amostra, e corta
\textbf{variância}; o não paramétrico é para o modelo rígido demais para a verdade,
e corta \textbf{viés}. A Figura~3.8 do [AME] (§3.9) desenha a relação entre os dois.
\end{observacao}

%==============================================================================
\section{Estratégia de bases: séries ortogonais e \emph{splines}}
%==============================================================================
A primeira ideia para ganhar flexibilidade sem sair do arcabouço linear é
\textbf{ampliar o conjunto de atributos} com funções $\phi_1,\dots,\phi_I$ das
covariáveis e ajustar
\[
  g(x) = \sum_{j=1}^{I} \beta_j \,\phi_j(x),
\]
que é \emph{linear nos $\beta_j$} --- estimável por mínimos quadrados (Aula~02).
Aqui a flexibilidade é controlada por $I$ (número de funções da base).

\subsection{Séries ortogonais}
Usa-se uma base ortonormal de $L^2$ (Fourier, polinômios ortogonais, \emph{wavelets}).
Aproxima-se $r$ por suas primeiras $I$ componentes na base; $I$ é o \emph{tuning
parameter} (mais termos $\Rightarrow$ mais flexível). É a ideia por trás de
``regressão polinomial'' vista na Aula~01 (a base dos monômios $1,x,x^2,\dots$).

\subsection{\emph{Splines}}
\emph{Splines} são funções polinomiais por partes, ``coladas'' suavemente em
pontos chamados \textbf{nós} (\emph{knots}) $t_1,\dots,t_p$. Uma base conveniente
para \emph{splines} de grau $k$ é
\[
  f_1(x)=1,\ f_2(x)=x,\ \dots,\ f_{k+1}(x)=x^k,\qquad
  f_{k+1+j}(x)=(x-t_j)_+^{k},\ j=1,\dots,p,
\]
onde $(u)_+=\max(u,0)$. Com isso $I=k+1+p$, e novamente estimam-se os
coeficientes por mínimos quadrados. Variantes importantes: \emph{B-splines}
(base numericamente estável), \emph{natural splines} (lineares fora dos nós
extremos, reduzindo a variância na fronteira) e \emph{smoothing splines} (penalizam
a curvatura --- caso particular de RKHS, §4.6).

\begin{emsala}
Mostre graficamente: poucos nós $\to$ curva rígida (subajuste); muitos nós $\to$
curva ``nervosa'' (superajuste). O número/posição dos nós é o que controla a
complexidade. Conecte: ``isto é a Aula~01 de novo''.
\end{emsala}

%==============================================================================
\section{$k$ vizinhos mais próximos (KNN)}
%==============================================================================
O KNN é talvez o método não paramétrico mais intuitivo: para prever em $\x$, faça
a \emph{média das respostas dos $k$ pontos de treino mais próximos de $\x$}.

\begin{definicao}[Regressão KNN]
Seja $\mathcal{N}_{\x}$ o conjunto dos índices das $k$ observações de treino mais
próximas de $\x$ (segundo uma distância $p$, em geral euclidiana):
$\mathcal{N}_{\x}=\{i : d(\X_i,\x)\le d^{k}_{\x}\}$, onde $d^k_{\x}$ é a distância
do $k$-ésimo vizinho. Então
\begin{equation}\label{eq:knn}
  \widehat{r}(\x) = \frac{1}{k}\sum_{i\in\mathcal{N}_{\x}} Y_i .
\end{equation}
\end{definicao}

O parâmetro $k$ controla a flexibilidade:
\begin{itemize}[leftmargin=1.4cm]
  \item \textbf{$k$ pequeno} (ex.: $k=1$): vizinhança minúscula, segue cada ponto
        $\Rightarrow$ \emph{viés baixo, variância alta} (curva ``nervosa'',
        superajuste).
  \item \textbf{$k$ grande}: média sobre muitos pontos; quando $k\to n$, $\rhat$
        tende à média global (constante) $\Rightarrow$ \emph{viés alto, variância
        baixa} (subajuste).
\end{itemize}
Escolhe-se $k$ por validação cruzada.

\begin{atencao}
KNN exige uma noção de distância e, portanto, \textbf{depende fortemente da escala}
das covariáveis: uma variável em unidade grande domina a distância. \emph{Padronize}
antes (Aula~07). Além disso, sofre da \emph{maldição da dimensionalidade}: em $p$
alto, ``vizinhos'' deixam de ser próximos (tema da Aula~05).
\end{atencao}

%==============================================================================
\section{Uma visão unificadora: suavizadores lineares}
%==============================================================================
KNN, \emph{séries} e \emph{splines} são todos
\textbf{suavizadores lineares}: a predição é uma combinação linear das respostas,
\[
  \widehat{r}(\x) = \sum_{i=1}^n \ell_i(\x)\,Y_i,
\]
com pesos $\ell_i(\x)$ que dependem só das covariáveis (não de $\bm{Y}$). Para o
KNN, $\ell_i(\x)=\tfrac1k\1{i\in\mathcal{N}_{\x}}$; para uma base, $\ell_i(\x)$ vem
do ajuste de mínimos quadrados na base.
Essa estrutura comum explica por que todos têm o mesmo dilema (o botão de
flexibilidade $\leftrightarrow$ viés--variância) e por que valem atalhos como o do
LOOCV (Aula~03, Prop.\ do \emph{leverage}).

%==============================================================================
\section*{Resumo}
%==============================================================================
\begin{itemize}[leftmargin=1.4cm]
  \item Métodos não paramétricos não fixam a forma de $r$; reduzem viés e aumentam
        variância em relação aos paramétricos.
  \item \emph{Bases} (séries, \emph{splines}): linearizam o problema; flexibilidade
        via número de termos/nós.
  \item \textbf{KNN} \eqref{eq:knn}: média dos $k$ vizinhos; $k$ controla
        viés--variância.
  \item Todos são \emph{suavizadores lineares}; \emph{padronize} as covariáveis e
        cuidado com a dimensão (Aula~05).
\end{itemize}

\paragraph{Leitura recomendada.} [AME] §4.1--4.3 (séries, \emph{splines}, KNN);
[ISLP] §3.5 (comparação regressão linear $\times$ KNN) e §7.1--7.5
(\emph{Moving Beyond Linearity}: \emph{splines}).

\end{document}
